\documentclass[notes=margin]{texign-article}

\subject{Seminar de istorie modernă}
\title{Școală, stat și națiune (1864–1918)}
\subtitle{Cum a devenit școala primară un proiect de stat}
\author{István}
\date{noiembrie 2026}
\addbibresource{surse.bib}

\begin{document}
\maketitle

\begin{abstract}
Lucrarea urmărește drumul de la Legea instrucțiunii din 1864 la reformele lui Spiru Haret și limitele proiectului școlar până la Marea Unire.
\end{abstract}

\section{Legea instrucțiunii publice}
Legea din 1864 a introdus învățământul primar obligatoriu și gratuit\cite[art.~31]{lege1864}, dar aplicarea ei a depins de resursele fiecărei comune. Rapoartele revizorilor școlari arată o realitate inegală: în unele sate școala funcționa într-o casă închiriată, în altele lipsea cu totul\cite[112]{murgescu2010}.

\begin{definition}[Sursă primară]
Document produs în perioada studiată, de un participant sau martor al evenimentelor.
\end{definition}

Statul își asumase o misiune pe care nu o putea încă finanța. Generația care a urmat a transformat școala într-un instrument al construcției naționale.

\image{revizor-1872.jpg}{Raport de revizor școlar, 1872.}

\begin{source}{Legea instrucțiunii, 1864, art.~31}
„Instrucțiunea primară este obligatorie și gratuită.”
\end{source}

\subsection{Învățământul primar}
Programele, manualele și formarea învățătorilor au devenit subiecte de dezbatere publică\cite[45--48]{iorga1928}. Revizorii au raportat sistematic lipsurile\cite{raport1872}. Ideea apare și în presa de specialitate\cite[14]{popescu1929}. Aceeași sursă insistă asupra salariilor\cite[16]{popescu1929}. Aceleași concluzii apar și la Murgescu\cite[114]{murgescu2010}.\note{Datele provin din anuarele statistice ale Ministerului Instrucțiunii, 1890–1914.}

\begin{quotation}[Nicolae Iorga, \emph{Istoria învățământului românesc}, 1928]
Școala satului nu poate fi mai bună decât învățătorul ei.
\end{quotation}

\begin{table}[htbp]
\caption{Școli primare rurale și elevi înscriși, 1864–1914}
\begin{tabular*}{\linewidth}{@{\extracolsep{\fill}}lrrr@{}}
\toprule
\head{An} & \head{Școli} & \head{Elevi} & \head{Elevi / școală}\\
\midrule
1864 & 2\,036 & 61\,000 & 30\\
1890 & 2\,654 & 124\,000 & 47\\
1914 & 4\,730 & 553\,000 & 117\\
\bottomrule
\end{tabular*}
\end{table}

\begin{warning}
Nu confundați data adoptării legii cu data intrării ei în vigoare.
\end{warning}

\section{Reformele lui Spiru Haret}
Spiru Haret a legat explicit educația de dezvoltarea satului.\note{Spiru Haret (1851–1912), ministru al instrucțiunii în trei rânduri.} Între 1890 și 1914 numărul elevilor s-a mai mult decât dublat, iar $\Delta = \frac{553\,000}{124\,000} \approx 4{,}5$.

\image{portret-haret.jpg}{Spiru Haret, cca.~1905.}

\begin{imagegrid}{Clădiri școlare și extinderea rețelei, 1890–1914.}
\gridimage{scoala-valcea.jpg}{Școală rurală, Vâlcea}
\gridimage{scoala-iasi.jpg}{Școală urbană, Iași}
\gridimage{harta-1890.pdf}{Rețeaua școlară, 1890}
\gridimage{harta-1914.pdf}{Rețeaua școlară, 1914}
\end{imagegrid}

\begin{summary}
\begin{itemize}
\item Învățământ primar obligatoriu din 1864.
\item Rețeaua școlară se dublează până în 1914.
\end{itemize}
\end{summary}

Pentru o perspectivă comparativă, vezi \textcite[cap.~5]{hitchins1994}.

\printsources
\end{document}
